```latex
% ============================================================
% 3.4.4 OVERALL RESIDUAL RISK EVALUATION
% ============================================================

\hypertarget{p3-4-4}{}

\subsection{Overall Residual Risk Evaluation}

\subsubsection{General}

The overall residual risk evaluation for the \textbf{BN-100 LUX} Infant
Radiant Warmer is addressed through the manufacturer's risk management
documentation.

Residual risk refers to the risk remaining after the documented
risk-control measures have been implemented.

The overall assessment considers the relationship between the identified
hazards, the implemented risk-control measures and the risks remaining
after application of those measures.

\subsubsection{Risk Management Basis}

The available manufacturer risk management document is identified as:

\begin{itemize}

    \item \textbf{Risk Management Report (BN-100/BN-100A)};

    \item Document No.: \textbf{YF-7.3-24-02 (w)};

    \item Version: \textbf{A/0}.

\end{itemize}

The title and model designation of the source document are retained as
provided by the manufacturer. The medical device covered by the present
dossier remains the \textbf{BN-100 LUX}.

\subsubsection{Residual Risks Following Risk Control}

The manufacturer has implemented documented risk-control measures addressing
the identified electrical, thermal, mechanical, monitoring and alarm-related
hazards.

The principal documented controls include:

\begin{itemize}

    \item electrical protective components;

    \item temperature monitoring;

    \item temperature control;

    \item over-temperature protection;

    \item heating cut-off;

    \item audible and visual alarms;

    \item sensor fault detection;

    \item power failure indication;

    \item mechanical operating limits;

    \item manufacturing and assembly controls.

\end{itemize}

After implementation of these measures, residual risks may remain as an
inherent consequence of the intended use and operating principle of the
device.

\subsubsection{Thermal Residual Risk}

The \textbf{BN-100 LUX} provides radiant heat directly to the newborn or
infant.

Consequently, thermal risk remains an inherent consideration even when
temperature control, monitoring and over-temperature protection are
provided.

The documented temperature-control and alarm functions constitute the
manufacturer's measures for controlling this risk.

The residual thermal risk is therefore considered within the manufacturer's
overall risk management process.

\subsubsection{Electrical Residual Risk}

Electrical hazards may remain as residual risks associated with the use of
medical electrical equipment.

The manufacturer addresses these risks through the documented electrical
design, protective components and safety-related functions.

The remaining electrical risk is considered within the manufacturer's
overall risk management framework.

\subsubsection{Mechanical Residual Risk}

Mechanical risks may remain during operation, positioning, adjustment or
loading of the patient support structure and associated components.

The manufacturer specifies mechanical limits and operating conditions for
these components.

Compliance with the specified operating limits and the manufacturer's
instructions forms part of the control of the associated residual risks.

\subsubsection{Temperature Monitoring Residual Risk}

Failure or incorrect operation of the skin temperature monitoring system may
result in residual risk.

The manufacturer addresses documented sensor failure conditions through
specific alarm functions, including short-circuit and open-circuit
detection.

The remaining risk associated with temperature monitoring is considered
within the overall manufacturer's risk management process.

\subsubsection{Alarm-Related Residual Risk}

Alarm functions are intended to alert the operator to specified abnormal
conditions.

A residual risk remains if an abnormal condition is not detected, if the
alarm is not appropriately perceived by the operator, or if the device is
not operated in accordance with the manufacturer's instructions.

The manufacturer documentation addresses alarm-related risks through the
documented alarm functions and operating instructions.

\subsubsection{Information for Safety}

Information provided in the manufacturer's user documentation forms part of
the overall risk-control strategy.

The operating instructions provide information concerning:

\begin{itemize}

    \item intended operation;

    \item temperature-control modes;

    \item alarm conditions;

    \item environmental operating conditions;

    \item storage conditions;

    \item applicable operating limitations.

\end{itemize}

Correct use of the device in accordance with the manufacturer's instructions
is therefore part of the management of residual risk.

\subsubsection{Overall Residual Risk}

The overall residual risk is evaluated by considering the risks remaining
after implementation of the documented risk-control measures.

The available manufacturer documentation provides the basis for this
evaluation through the identified hazards, risk-control measures and
corresponding safety functions.

No independent numerical residual-risk value or new risk-ranking result is
introduced in this section where such information is not explicitly
available in the manufacturer's documentation.

\subsubsection{Benefit--Risk Consideration}

The \textbf{BN-100 LUX} is intended to provide controlled radiant thermal
support to newborn and infant patients requiring thermal care.

The intended clinical function of the device is considered together with
the risks identified and controlled through the manufacturer's risk
management process.

The overall benefit--risk consideration shall be understood in conjunction
with the manufacturer's intended-use information, risk management
documentation and applicable technical evidence.

\subsubsection{Documentation Limitation}

The present section summarizes the overall residual-risk considerations
supported by the available manufacturer documentation.

It does not replace the original manufacturer's risk management file and
does not introduce independent residual-risk scores, acceptance criteria or
benefit--risk conclusions that are not documented by the manufacturer.

Where the competent authority requires detailed residual-risk records, the
original Risk Management Report and its associated risk-analysis tables
should be included as supporting documentation.

\subsubsection{Documentary Basis}

This section is based primarily on:

\begin{itemize}

    \item Risk Management Report (BN-100/BN-100A), Document No.
    YF-7.3-24-02 (w), Version A/0;

    \item BN-100 / BN-100 LUX User Manual.

\end{itemize}
```
